While it is common to use MLTT as an internal language for
certain arguments about objects in its models,
this idea has not yet been fully implemented in a proof assistant.
The HoTTLean project \cite{hua2025, nawrocki2026, hottuf2026}
(\href{https://github.com/sinhp/HoTTLean}{github.com/sinhp/HoTTLean}) 
aims to bridge this gap by
designing a domain-specific language (DSL) for Martin-L\"of type theories (MLTT)
using Lean's DSL support,
formalising the general semantics of MLTT in Lean,
and supplying a specific model
-- namely the Hofmann--Streicher groupoid model \cite{hofmann1995} --
as a test case.
Overall, this allows users to write synthetic proofs in the DSL,
and use the interpretation to produce constructions pertaining to
groupoids as defined in Mathlib \cite{mathlib2020},
the standard library for mathematics in Lean.
In short, this would make automated internal reasoning in the category of groupoids a reality.

Here, we will focus on the formalisation of the groupoid model,
with a finite hierarchy of universes,
$\Si$-types, $\Pi$-types, and $\Id$-types.

The process of formalisation involves a great deal of experimentation.
It is often difficult to ascertain what approach is best until it has been tried.
For the construction of a model of type theory,
there are many choices to be made.
For example,
just the definition of ``a model of MLTT'' can be taken to mean
a natural model \cite{awodey2018},
or a category with families (CwF) \cite{clairambault2011},
or a comprehension category \cite{jacobs1993};
the list goes on.
After various experiments,
we decided on an approach that was (1) ``internal'',
and (2) ``unalgebraic''.

% The word ``internal'' here refers to replacing constructions on an external classifier
% of all type families (in presheaves) with 
% constructions on classifiers of small type families (in the category itself).

By ``internal'', we mean that instead of working with a classifier of all types $\Ty : \CC^\op \to \Set$
(as is typically required for categories with families and natural models),
we work with a hierarchy of universes in the category of contexts
$U_0 \to U_1 \to \cdots \in \CC$
(i.e., a hierarchy of representable presheaves $\Ty_0 \to \Ty_1 \to \cdots : \CC^\op \to \Set$
classifying types of a certain size).
% such that the universes are defined polymorphically over the universe levels of Lean.
Working internally meant that the user supplying a model need not make separate constructions
for the classifier of all types and the internal universes classifying small types.
It also simplified the universe level calculus,
as presheaves involve introducing extra universe levels and size constraints therein.
Furthermore, this approach seems to be more uniform across different models of HoTT.
Usually, an ad hoc description of simplical and cubical models is given in a presheaf category,
to which we would not want to apply the theory of natural models or CwFs.
In short, we are satisfied with the internal approach.

On the other hand,
the ``unalgebraic'' construction of the groupoid model is
not entirely satisfactory.
(Recall that we refer to the natural models style of semantics as \emph{algebraic}
(as in \emph{algebraic type theory} \cite{awodey2025}),
and use the word \emph{unalgebraic} in contrast
to describe semantics with explicit coherence conditions, like a CwF.)
We make the case that an alternative ``algebraic'' approach,
presented in \Cref{sec:StrSem},
would improve upon this first attempt
(in terms of the metrics for formalisation).
We discuss this in \Cref{sec:algebraic-vs-elementary}.

In the following, we present general unalgebraic models
and the groupoid model constructed as an unalgebraic model,
as a companion to the formalisation.
We also present algebraic models and a roadmap towards formalisation
using algebraic models.
The formalisation repository can be found at
\href{https://github.com/sinhp/HoTTLean/tree/HoTTUF}{\textsf{github.com/sinhp/HoTTLean}}.
We also include links to the
formalised counterparts of definitions,
labeled using the symbol \done{https://github.com/sinhp/HoTTLean/tree/HoTTUF}.

\subsection*{Related Work}
Du \cite{du2025} formalised in Lean a simplicial model
of type theory as a contextual category with $\Pi$-types,
a single universe, and without other type formers.
Sozeau and Tabareau \cite{sozeau2014} worked towards a Rocq formalisation
of the groupoid model,
but this has not been completed, to our knowledge.
The groupoid model is formalised as an example comprehension category
in the Unimath project \cite{2024grayson},
also without type formers.

\section{Unalgebraic models}\label{sec:UnstrSem}
We fix a category $\CC$,
and call objects of $\CC$ \emph{contexts} and
morphisms of $\CC$ \emph{substitutions}. 

\begin{definition}[\done{https://github.com/sinhp/HoTTLean/blob/HoTTUF/HoTTLean/Model/Unstructured/UnstructuredUniverse.lean\#L19}]
  \label{def:elementary-universe}
  A \emph{universe}
  in $\CC$ consists of a morphism
  $\uu : \Ulow \to \U$
  and a chosen pullback $\Gamma . A$
  for each object $\Gamma$ and morphism $A : \Gamma \to \U$.
  \[\begin{tikzcd}
	{\Gamma.A} & \Ulow \\
	\Gamma & \U
	\arrow["\var_A", from=1-1, to=1-2]
	\arrow["d_A", swap, from=1-1, to=2-1]
	\arrow["\lrcorner"{anchor=center, pos=0.125}, draw=none, from=1-1, to=2-2]
	\arrow["\uu", from=1-2, to=2-2]
	\arrow["A"', from=2-1, to=2-2]
  \end{tikzcd}\]
  We will view $A : \Gamma \to \U$,
  as a type in context $\Gamma$,
  and view the chosen pullback $\Gamma . A$
  as context extension $d_A : \Gamma . A \to \Gamma$.
  We also view a morphism $a : \Gamma \to \Ulow$
  satisfying $\uu \circ a = A$ as a term of type $A$.
\end{definition}

One can view $\uu : \Ulow \to \U$ as a context extension
$d_X : (X : \U.\,x : X) \to (X : \U)$,
by taking the type $X : \U \to \U$ to be the identity
$\id_\U : \U \to \U$.

\begin{definition}[\done{https://github.com/sinhp/HoTTLean/blob/HoTTUF/HoTTLean/Model/Unstructured/UnstructuredUniverse.lean\#L98}]
  For $\sigma : \Delta \to \Gamma$,
  a type $A : \Gamma \to \U$ and a term $a : \Delta \to \Ulow$
  such that $\uu \circ a = A \circ \sigma$,
  we denote the map into the context extension
  using the notation
  $\sigma . a : \Delta \to \Gamma . A$,
  induced by the pullback
\[\begin{tikzcd}
	\Delta \\
	& {\Gamma.A} & \Ulow \\
	& \Gamma & \U
	\arrow["{\sigma.a}"{description}, dashed, from=1-1, to=2-2]
	\arrow["a", bend left, from=1-1, to=2-3]
	\arrow["\sigma"', bend right, from=1-1, to=3-2]
	\arrow[from=2-2, to=2-3]
	\arrow[from=2-2, to=3-2]
	\arrow["\lrcorner"{anchor=center, pos=0.125}, draw=none, from=2-2, to=3-3]
	\arrow["\uu", from=2-3, to=3-3]
	\arrow["A"', from=3-2, to=3-3]
\end{tikzcd}\]
\end{definition}

\begin{definition}[$\Unit$-types] \label{def:elementary-unit}
  We say that a universe $\uu : \Ulow \to \U$ admits
  \emph{unalgebraic $\Unit$-types} when we have the following.
  \begin{enumerate}
  \item
    For any context $\Gamma$, there is a type $\Unit_\Gamma : \Gamma \to \U$.
  \item
    The construction $\Unit_\Gamma$ is stable under substitution,
    meaning that for any $\sigma : \Delta \to \Gamma$ we have
    \[
      \Unit_\Gamma \circ \sigma = \Unit_\Delta
    \]
  \item
    For any context $\Gamma$, there is a unique map $\unit_\Gamma : \Gamma \to \Ulow$
    satisfying $\uu \circ \unit_\Gamma = \Unit_\Gamma$.
  \end{enumerate}
\end{definition}

\begin{definition}[$\Sigma$-types, \done{https://github.com/sinhp/HoTTLean/blob/HoTTUF/HoTTLean/Model/Unstructured/UnstructuredUniverse.lean\#L219}]
  \label{def:elementary-sigma}
  We will say $\uu : \Ulow \to \U$ admits
  \emph{unalgebraic $\Sigma$-types}
  when we have the following.
  \begin{enumerate}
  \item 
    For any context $\Gamma$, type $A : \Gamma \to \U$,
    and dependent type $B : \Gamma . A \to \U$
    there is a type $\Sigma_A B : \Gamma \to \U$.
  \item
    The construction $\Sigma_A B$ is stable under substitution,
    meaning that for any $\sigma : \Delta \to \Gamma$ we have
    \[
      \Sigma_{A \circ \sigma} (B \circ \tilde{\sigma}) =
      \Sigma_{A} B \circ \sigma
    \]
    where $\tilde{\sigma}$ is the same as in \Cref{def:elementary-pi}.
  \item
    Under the assumptions of (1),
    for any $a : \Gamma \to \Ulow$ and $b:\Gamma\to \Ulow$
    such that $\uu \circ a = A$ and $\uu \circ b = B \circ \id_\Gamma . a $,
    there is a map $\pair (a,b) : \Gamma \to \Ulow$
    satisfying $\uu \circ \pair (a,b)  = \Sigma_A B$. 
    \[
\begin{tikzcd}
	\Gamma \\
	& {\Gamma . A} & \Ulow \\
	& \Gamma & \U
	\arrow["{\id_\Gamma . a}"{description}, from=1-1, to=2-2]
	\arrow["a", bend left, from=1-1, to=2-3]
	\arrow["{\id_\Gamma}"', bend right, from=1-1, to=3-2]
	\arrow[from=2-2, to=2-3]
	\arrow[from=2-2, to=3-2]
	\arrow["\uu", from=2-3, to=3-3]
	\arrow["A"', from=3-2, to=3-3]
\end{tikzcd}
    \]
  \item
    The construction $\pair (a, b)$ is stable under substitution.
    For any $\sigma : \Delta \to \Gamma$
    \[
      \pair (a \circ \sigma,b \circ \sigma) =
       \pair(a , b) \circ \sigma
    \]
  \item
    Under the assumptions of (1),
    for any $s : \Gamma \to \Ulow$
    such that $ \uu \circ s = \Sigma_A B$,
    there are maps $\fst s : \Gamma  \to \Ulow$
    and $\snd s : \Gamma  \to \Ulow$
    satisfying $\uu \circ \fst\; s = A$
    and $\uu \circ \snd \; s = B \circ (\id_\Gamma . \fst\; s)$.
  \item ($\beta$-rules) The operations 
    $\fst$ and $\snd$ deconstruct a term formed using $\pair$.
    \[ \fst \;(\pair(a, b)) = a \quad \text{ and } 
    \quad \snd \; (\pair (a, b)) = b\]
  \item ($\eta$-rule) Conversely, the operation $\pair$
    reconstructs the term split up by $\fst$ and $\snd$.
  \[\pair (\fst \;s, \snd \;s) = s\]
  \end{enumerate}
\end{definition}

In the presence of sufficient exponentiability conditions --
for example, if there is an embient $\pi$-clan (\Cref{def:pi-preclan}) or
$\CC$ is lex and the universe is an exponentiable morphism
(see \Cref{sec:exponentiable} and \Cref{sec:polynomials}) --
then unalgebraic $\Sigma$-types on $\uu$
can be reformulated algebraically using polynomials.
We provide such a reformulation in \Cref{def:alg-sig-type-mltt}.

% \[
%     \begin{tikzcd}
%       Q & \Ulow \\
%       {P_\uu \U} & \U
%       \arrow["{\pair}", from=1-1, to=1-2]
%       \arrow["{{\uu \pcomp \uu}}"', from=1-1, to=2-1]
%       \arrow["\lrcorner"{anchor=center, pos=0.125}, draw=none, from=1-1, to=2-2]
%       \arrow["\uu", from=1-2, to=2-2]
%       \arrow["{\Sigma}"', from=2-1, to=2-2]
%     \end{tikzcd}
% \]

For reasons of generality that will become important in \Cref{sec:cat-as-types},
we will define a version of $\Pi$-types that is slightly more general than what is needed in this section,
involving two universes rather than just one.
(In fact, the formalisation is even more general,
involving three universes. See \Cref{rmk:three-universes-generality}.)
\begin{definition}[$\Pi$-types, \done{https://github.com/sinhp/HoTTLean/blob/HoTTUF/HoTTLean/Model/Unstructured/UnstructuredUniverse.lean\#L332}]
  \label{def:elementary-pi}
  Let $\uu : \Ulow \to \U$ and $\vv : \Vlow \to \V$ be two universes.
  We will say that $\uu$ admits $\vv$-indexed \emph{unalgebraic $\Pi$-types}
  when we have the following.
  \begin{enumerate}
  \item 
    For any context $\Gamma$, type $A : \Gamma \to \V$,
    and dependent type $B : \Gamma . A \to \U$
    there is a type \[\Pi_A B : \Gamma \to \U\]
  \item
    The construction $\Pi_A B$ is stable under substitution,
    meaning that for any $\sigma : \Delta \to \Gamma$ we have
    \[
      \Pi_{A \circ \sigma} (B \circ \tilde{\sigma}) =
      B \circ \sigma
    \]
    where $\tilde{\sigma} := (\sigma \circ d_{A \circ \sigma}) . \var_{A \circ \sigma}$ as in
    \[
    \begin{tikzcd}
      {\Delta . {(A \circ \sigma)}} & {\Gamma . A} & \Vlow \\
      \Delta & \Gamma & \V
      \arrow["{{\tilde{\sigma}}}", dashed, from=1-1, to=1-2]
      \arrow["{\var_{A \circ \sigma}}", bend left, from=1-1, to=1-3]
      \arrow["{d_{A \circ \sigma}}", from=1-1, to=2-1]
      \arrow["{\var_A}", from=1-2, to=1-3]
      \arrow["{d_A}", from=1-2, to=2-2]
      \arrow["\lrcorner"{anchor=center, pos=0.125}, draw=none, from=1-2, to=2-3]
      \arrow[from=1-3, to=2-3]
      \arrow["\sigma"', from=2-1, to=2-2]
      \arrow["A"', from=2-2, to=2-3]
    \end{tikzcd}
    \]
  \item
    Under the assumptions of (1),
    for any $b : \Gamma . A \to \Ulow$
    such that $\uu \circ b = B$,
    there is a map $\lam  b : \Gamma \to \Ulow$
    satisfying $\uu \circ \lam \;b  = \Pi_A B$.
  \item
    The construction $\lam  b$ is stable under substitution.
    For any $\sigma : \Delta \to \Gamma$
    \[
      \lam (b \circ \tilde{\sigma}) =
      \lam  b \circ \sigma
    \]
  \item
    Under the assumptions of (1),
    for any $f : \Gamma \to \Ulow$
    such that $\uu \circ f = \Pi_A B$,
    there is a map $\unlam  f : \Gamma . A \to \Ulow$
    satisfying $\uu \circ \unlam  f  = B$.
  \item ($\beta$- and $\eta$-rules) The operations 
    $\lam$ and $\unlam$ are inverses of one another.
    \[ \unlam (\lam  b) = b \quad \text{ and } 
    \quad \lam (\unlam  f) = f\]
  \end{enumerate}
  If $\uu$ admits unalgebraic $\uu$-indexed $\Pi$-types,
  then we will simply say that $\uu$ admits
  \emph{unalgebraic $\Pi$-types}.
\end{definition}
Again, when polynomials are available,
unalgebraic $\Pi$-types on the universe can be 
reformulated algebraically as a pullback square (\Cref{def:alg-pi-type-mltt}).

% In the presence of sufficient exponentiability conditions,
% for example if there is an ambient $\pi$-clan or
% $\CC$ is lex and the universe is an exponentiable morphism
% (see \Cref{sec:exponentiable} and \Cref{sec:polynomials}),
% then an elementary $\Pi$-types on $\uu$
% can be reformulated algebraically using polynomials.
% We provide such this reformulation in \Cref{def:alg-pi-type-mltt}.

\begin{definition}[$\Id$-types, \done{https://github.com/sinhp/HoTTLean/blob/HoTTUF/HoTTLean/Model/Unstructured/UnstructuredUniverse.lean\#L439}]
  We say that $\uu : \Ulow \to \U$
  admits \emph{unalgebraic $\Id$-types}
  when we have the following.
  \begin{enumerate}
  \item 
    For any context $\Gamma$, type $A : \Gamma \to \U$,
    and terms $a_0, a_1 : \Gamma \to \Ulow$ such that
    $\uu \circ a_i = A$,
    there is a type $\Id_A (a_0,a_1) : \Gamma \to \U$.
  \item
    The construction $\Id_A (a_0,a_1)$ is stable under substitution.
  \item
    For any context $\Gamma$, type $A : \Gamma \to \U$,
    and term $a : \Gamma \to \Ulow$ such that
    $\uu \circ a = A$,
    there is a term $\refl a : \Gamma \to \Ulow$
    satisfying
    $\uu \circ \refl a  = \Id_A (a,a)$.
  \item The construction $\refl$ is stable under substitution.
  \item
    Under the assumptions of (3)
    one can construct the context $\Gamma . (x : A) . \Id_A(a,x)$
    by taking context extension on $\Gamma . A$ for the type
    $\Id_A(a,x)$ given by (1),
    and the substitution
    \[\rho_a := \id_\Gamma . a . \refl a : \Gamma \to \Gamma . (x : A) . \Id_A(a,x)\]
    For any \emph{motive} $C : \Gamma .(x : A) . \Id_A(a,x) \to \U$
    -- that means $C$ is parametrised by a family of paths with a fixed starting point $a$ --
    and map $r : \Gamma \to \Ulow$ satisfying
    \[ \uu \circ r = C \circ \rho_a \]
    there is a term
    $j (C, r) : \Gamma . (x : A) . \Id_A(a,x) \to \Ulow$
    satisfying $\uu \circ j (C, r) = C$ and
    $j (C, r) \circ \rho_a  = r$
    \[\begin{tikzcd}
	   \Gamma & \Ulow \\
	   {\Gamma . (x : A) . \Id_A(a,x)} & \U
	   \arrow["{r}", from=1-1, to=1-2]
	   \arrow["{\rho_a}"', from=1-1, to=2-1]
	   \arrow["\uu", from=1-2, to=2-2]
	   \arrow["{j(C,r)}"{description}, dashed, from=2-1, to=1-2]
	   \arrow["C"', from=2-1, to=2-2]
    \end{tikzcd}\]
  \item The construction $j$ is stable under substitution.
  \end{enumerate}
\end{definition}

\Cref{def:algebraic-id-types} provides the algebraic version of identity type structure.

\begin{remark}\label{rmk:three-universes-generality}
% In the previous definitions,
% we have exclusively worked with a single universe $\uu$,
% for the sake of simplicity.
% However,
The definitions for $\Sigma$-types (and similarly $\Pi$-types)
% $\Sigma$-types in Lean
have $A$ in a universe $\uu_0$,
$B$ in a second universe $\uu_1$,
and $\Sigma_A B$ in a third universe $\uu_2$.
% A similar generalisation can be made for path types.
The formalisation also makes $\Id$-types \emph{large eliminating},
meaning that the motive for identity elimination $C$
need not live in the same universe as the type $A$ on which
identities are taken.
The more general definitions can be found in the formalisation.
\end{remark}

The following defines
non-cumulative universe lifts
``\`a la Coquand'' \cite{gratzer2020}.

\begin{definition}[\done{https://github.com/sinhp/HoTTLean/blob/HoTTUF/HoTTLean/Model/Unstructured/UHom.lean\#L23},
\done{https://github.com/sinhp/HoTTLean/blob/HoTTUF/HoTTLean/Model/Unstructured/UHom.lean\#L77}]
  A \emph{morphism of universes} $l : \uu_0 \to \uu_1$ 
  consists of a pair of maps $l_\U : \U_0 \to \U_1$
  and $l_\Ulow : \Ulow_0 \to \Ulow_1$ such that 
  the following square is a pullback
  \[\begin{tikzcd}
	{\Ulow_0} & {\Ulow_1} \\
	{\U_0} & {\U_1}
	\arrow["{{{l_\Ulow}}}", from=1-1, to=1-2]
	\arrow["{{{\uu_0}}}"', from=1-1, to=2-1]
	\arrow["\lrcorner"{anchor=center, pos=0.125}, draw=none, from=1-1, to=2-2]
	\arrow["{\uu_1}", from=1-2, to=2-2]
	\arrow["{{{l_\U}}}"', from=2-1, to=2-2]
\end{tikzcd}\]
  Suppose $\CC$ has a terminal object $1$.
  A \emph{universe lift} consists of
  a morphism of universes $l : \uu_0 \to \uu_1$
  and a pair of maps $\ulcorner U_0 \urcorner : 1 \to \U_1$
  and $a : \U_0 \to \Ulow_1$
  such that the following square is a pullback.
\[\begin{tikzcd}
	{\U_0} & {\Ulow_1} \\
	1 & {\U_1}
	\arrow["{{a}}", from=1-1, to=1-2]
	\arrow[from=1-1, to=2-1]
	\arrow["\lrcorner"{anchor=center, pos=0.125}, draw=none, from=1-1, to=2-2]
	\arrow["{{{\uu_1}}}", from=1-2, to=2-2]
	\arrow["{{{\ulcorner U_0 \urcorner}}}"', from=2-1, to=2-2]
\end{tikzcd}\]
\end{definition}

\subsection*{Unalgebraic path types}

Path types will be covered in more detail in \Cref{sec:path-types},
in an algebraic style.
For HoTTLean, an unalgebraic formulation of path types
was used in the HoTTLean formalisation
to construct identity types in the groupoid model.
In this section, we briefly introduce the relevant definitions
of unalgebraic path types,
and leave a more detailed discussion to \Cref{sec:path-types}.

\begin{definition}[\done{https://github.com/sinhp/HoTTLean/blob/HoTTUF/HoTTLean/Model/Unstructured/Hurewicz.lean\#L22}]
  \label{def:cylinder-structure}
  A cylinder structure on a category $\CC$
  consists of
  \begin{itemize}
    \item an endofunctor ${\cyl} : \CC \to \CC$, called the \emph{cylinder}
    \item and natural transformations
      \[\delta_0, \delta_1 : \id_\CC \to {\cyl}
      \quad \quad
      \pi : {\cyl} \to \id_\CC
      \quad \quad 
      \psi : \cyl \circ \cyl \to \cyl \circ \cyl\]
    \item satisfying the following equations
    \[
    \begin{tikzcd}
      \id & {\cyl} & \id & {{\cyl} \circ {\cyl}} & {{\cyl} \circ {\cyl}} \\
      & \id &&& {{\cyl} \circ {\cyl}}
      \arrow["{\delta_0}", from=1-1, to=1-2]
      \arrow[equals, from=1-1, to=2-2]
      \arrow["\pi"', from=1-2, to=2-2]
      \arrow["{\delta_0}"', from=1-3, to=1-2]
      \arrow[equals, from=1-3, to=2-2]
      \arrow["\psi", from=1-4, to=1-5]
      \arrow[equals, from=1-4, to=2-5]
      \arrow["\psi", from=1-5, to=2-5]
    \end{tikzcd}
    \]
    \[ \begin{tikzcd}
      {\cyl} & {\cyl} & {\cyl} & {\cyl} & {{\cyl} \circ {\cyl}} & {{\cyl} \circ {\cyl}} \\
        {{\cyl} \circ {\cyl}} & {{\cyl} \circ {\cyl}} & {{\cyl} \circ {\cyl}} & {{\cyl} \circ {\cyl}} & {\cyl} & {\cyl}
        \arrow[equals, from=1-1, to=1-2]
        \arrow["{\delta_0 {\cyl}}"', from=1-1, to=2-1]
        \arrow["{{\cyl} \delta_0}", from=1-2, to=2-2]
        \arrow[equals, from=1-3, to=1-4]
        \arrow["{\delta_1 {\cyl}}"', from=1-3, to=2-3]
        \arrow["{{\cyl} \delta_1}", from=1-4, to=2-4]
        \arrow["\psi", from=1-5, to=1-6]
        \arrow["{\pi {\cyl}}"', from=1-5, to=2-5]
        \arrow["{{\cyl} \pi}", from=1-6, to=2-6]
        \arrow["\psi", from=2-1, to=2-2]
        \arrow["\psi", from=2-3, to=2-4]
        \arrow[equals, from=2-5, to=2-6]
      \end{tikzcd}\]
  \end{itemize}
\end{definition}

For example, a cylinder structure can be defined
using a bipointed object and binary products.
Suppose $\CC$ has a terminal object and $I$ is
a bipointed object $\delta_0, \delta_1 : 1 \to I$.
Then we can define the cylinder functor by $\cyl : X \mapsto X \times I$.

\begin{definition}[$\Path$-types, \done{https://github.com/sinhp/HoTTLean/blob/HoTTUF/HoTTLean/Model/Unstructured/Hurewicz.lean\#L238}]
  Suppose $\CC$ has a terminal object and
  a cylinder structure.
  A universe $\uu : \Ulow \to \U$ admits
  \emph{unalgebraic $\Path$-types}
  when we have the following.
  \begin{enumerate}
  \item 
    For any context $\Gamma$, type $A : \Gamma \to \U$,
    and terms $a_0, a_1 : \Gamma \to \Ulow$ such that
    $\uu \circ a_i  = A$,
    there is a type $\Path_A (a_0,a_1) : \Gamma \to \U$.
  \item
    The construction $\Path_A (a_0,a_1)$ is stable under substitution.
  \item
    For any context $\Gamma$, type $A : \Gamma \to \U$,
    and ``path'' (or ``homotopy'') $p : \cyl \Gamma \to \Ulow$ such that
    $\uu \circ p = A \circ \pi$,
    there is a term $\path  p : \Gamma \to \Ulow$
    satisfying
    \[\uu \circ \path  p  = {\Path}_A (p \circ \delta_0,p \circ \delta_1)\]
  \item The construction $\path$ is stable under substitution.
  \item
    For any context $\Gamma$, type $A : \Gamma \to \U$,
    a pair of terms $a_0, a_1 : \Gamma \to \Ulow$ such that
    $\uu \circ a_0 = \uu \circ a_1 = A$,
    and a term $p : \Gamma \to \Ulow$ such that
    $\uu \circ p = \Path_A (a_0,a_1)$,
    there is a term $\unpath  p : \cyl \Gamma \to \Ulow$
    satisfying
    \[\uu \circ \unpath  p = A \circ \pi
      \quad \quad
      \unpath p \circ \delta_0  = a_0
      \quad \quad
      \unpath p \circ \delta_1  = a_1
    \]
  \item Under the assumptions of (3),
    \[ \unpath (\path p) = p \]
  \item Under the assumptions of (5),
    \[ \path (\unpath p) = p \]
  \end{enumerate}
\end{definition}

Unalgebraic $\Path$-types are reformulated algebraically in
\Cref{def:mltt-alg-path-types-2}.

% The following is closely related to opfibrancy in $\Cat$
% (\Cref{def:split-opfibration}).
\begin{definition}[\done{https://github.com/sinhp/HoTTLean/blob/HoTTUF/HoTTLean/Model/Unstructured/Hurewicz.lean\#L148}]
  \label{def:hurewicz-defs}
  Suppose $\CC$ admits a cylinder structure.
  Let $f : A \to B$ be a morphism in $\CC$.
  A \emph{Hurewicz structure $(f,l)$} consists of
  \begin{itemize} 
    \item for every object $X$ in $\CC$,
    and every pair of maps $a : X \to A$
    and $p : \cyl X \to B$ such that
    $f \circ a = p \circ \delta_0$,
    a chosen diagonal filler $l_{(a,p)} : \cyl X \to A$,
    \[\begin{tikzcd}
        X & A \\
        {\cyl X} & B
        \arrow["a", from=1-1, to=1-2]
        \arrow["\delta_0"', from=1-1, to=2-1]
        \arrow["f", from=1-2, to=2-2]
        \arrow["{l_{(a,p)}}"{description}, dashed, from=2-1, to=1-2]
        \arrow["p"', from=2-1, to=2-2]
      \end{tikzcd} \]
    \item such that for any morphism $x : X' \to X$,
    the following coherence equation between lifts
    $l_{(a \circ x,p \circ \cyl x)} : \cyl X' \to A$
    and $l_{(a,p)} : \cyl X \to A$ holds.
    \[ l_{(a,p)} \circ \cyl x = l_{(a \circ x,p \circ \cyl x)} \]
  \end{itemize}

  We say that $(f,l)$ is \emph{normal} if
  for every $b : X \to B$, we have $l_{(a,b \circ \pi)} = a \circ \pi$. 
  \[\begin{tikzcd}
      X && A \\
      {\cyl X} & X & B
      \arrow["a", from=1-1, to=1-3]
      \arrow["{\delta_0}"', from=1-1, to=2-1]
      \arrow["f", from=1-3, to=2-3]
      \arrow["{l_{(a,b \circ \pi)}}"{description}, dashed, from=2-1, to=1-3]
      \arrow["\pi"', from=2-1, to=2-2]
      \arrow["a"{description}, from=2-2, to=1-3]
      \arrow["b"', from=2-2, to=2-3]
    \end{tikzcd}
    \]  
\end{definition}
When $\cyl = (-) \times I$ is given by product with an interval object,
a Hurewicz structure on $f : A \to B$ is the usual homotopy lifting
definition from \cite{awodey2026}.

\begin{theorem}[Identity elimination, \done{https://github.com/sinhp/HoTTLean/blob/HoTTUF/HoTTLean/Model/Unstructured/Hurewicz.lean\#L743}]
  \label{thm:id-elim-2}
  If $\uu : \Ulow \to \U$ admits unalgebraic $\Path$-types
  and a normal Hurewicz structure $(\uu,l)$,
  then $\uu$ also admits unalgebraic $\Id$-types.
\end{theorem}
A proof of \Cref{thm:id-elim-2} in the style of algebraic models is provided in \Cref{thm:id-elim}.

\section{The groupoid model}

\begin{definition}[Contexts \done{https://github.com/sinhp/HoTTLean/blob/31133dd5b25226ea897f8aa5e2e43b61392459eb/HoTTLean/Groupoids/Basic.lean\#L65}]
The category $\CC$ is taken to be the category of groupoids.
In Mathlib, there are two universe-level variables
associated with the category of groupoids,
one for the size of objects and one for the size of hom-sets,
both of which we (arbitrarily) take to be $4$.
\[ \CC := \Grpd.\{4,4\} \]
\end{definition}

\begin{definition}\label{def:split-isofibration}
  A cloven isofibration consists of a pair $(F,l)$,
  where $F : \DD \to \CC$ is a functor
  and $l$ provides for each object $x \in \DD$, each object $y \in \CC$,
  and each isomorphism $f : F x \to y$ in $\CC$,
  a chosen isomorphism $l_{(x,f)} : x \to y'$
  such that $F l_{(x,f)} = f$,
  called the lift of $f$.

  A cloven isofibration $(F : \DD \to \CC,l)$
  is \emph{normal} if 
  for any $x \in \DD$ the lift of the identity
  $\id_{F x} : F x \to F x$ is the identity
  \[l_{(x,\id_{F x})} = \id_x : x \to x\]
  In particular, the endpoint of the lift satisfies $(F x)' = x$.
  We call $(F,l)$ a \emph{normal isofibration} for short.

  A normal isofibration $(F : \DD \to \CC, l)$ is \emph{split} if
  for any $x \in \DD$, $y$ and $z$ in $\CC$ and
  isomorphisms $f : F x \to y$ and $g : y \to z$,
  the lift of the composition is the composition of the lifts
  \[l_{(x,g \circ f)} = l_{(x,g)} \circ l_{(y',f)} : x \to z' \]
  We call $(F,l)$ a \emph{split isofibration} for short.
\end{definition}

\begin{definition}[Universes and context extension \done{https://github.com/sinhp/HoTTLean/blob/31133dd5b25226ea897f8aa5e2e43b61392459eb/HoTTLean/Groupoids/UnstructuredModel.lean\#L26}]
\label{sec:universe-in-groupoid-model}
We provide a universe (as in \Cref{def:elementary-universe})
for each Lean universe level $\alpha < 4$ polymorphically.
$\uu_\alpha$ will be called the \emph{universe $\alpha$-small isofibrations}
(since each of its pullbacks admits a canonical isofibration structure):
\begin{itemize}
  \item $\U$ is the core of the category of ($\alpha$-small) groupoids
  \[\U_\alpha := \Core(\Grpd.\{\alpha,\alpha\})\]
  \item $\Ulow$ is the core of the category of ($\alpha$-small) pointed groupoids,
  \[\Ulow_\alpha := \Core(\PGrpd.\{\alpha,\alpha\})\]
  where the category of $\alpha$-small pointed groupoids is the
  Grothendieck construction on the forgetful functor from groupoids to categories
  $\Grpd.\{\alpha,\alpha\} \to \Cat.\{\alpha,\alpha\}$.
  \item The map $\uu$ is the image of the forgetful functor
  from pointed groupoids to groupoids $\PGrpd \to \Grpd$.
  \[ \uu_\alpha : \Ulow_\alpha \to \U_\alpha\]
  \item Pullbacks of each universe $\uu : \Ulow \to \U$
    are computed as Grothendieck constructions
  \[ \begin{tikzcd}
	{\int (i \circ A)} & {\Core(\PGrpd)} & \PGrpd \\
	\Gamma & {\Core(\Grpd)} & \Grpd
	\arrow[dashed, from=1-1, to=1-2]
	\arrow[from=1-1, to=2-1]
	\arrow["\lrcorner"{anchor=center, pos=0.125}, draw=none, from=1-1, to=2-2]
	\arrow[from=1-2, to=1-3]
	\arrow[from=1-2, to=2-2]
	\arrow["\lrcorner"{anchor=center, pos=0.125}, draw=none, from=1-2, to=2-3]
	\arrow[from=1-3, to=2-3]
	\arrow["A"', from=2-1, to=2-2]
	\arrow["i"', from=2-2, to=2-3]
  \end{tikzcd} \]
  where $i : \Core(\Grpd) \to \Grpd$ is the usual inclusion of the core category.
  We will use context extension notation for
  the pullback $\Gamma . A := {\int (i \circ A)}$.
  (We could alternatively have taken $\Ulow := \int i$ as the Grothendieck construction.)
\end{itemize}
\end{definition}

\begin{definition}[Universe lifts \done{https://github.com/sinhp/HoTTLean/blob/31133dd5b25226ea897f8aa5e2e43b61392459eb/HoTTLean/Groupoids/UHom.lean\#L33}]
Universes in the groupoid model are related by a chain of universe lifts
$\uu.\{\alpha\} \to \uu.\{\alpha+1\}$, for each $\alpha < 3$.
First, by lifting (Lean) types in universe $\alpha$ to universe $\alpha+1$,
we have a pullback square
\[
\begin{tikzcd}
	{\PGrpd.\{\alpha,\alpha\}} & {\PGrpd.\{\alpha+1,\alpha+1\}} \\
	{\Grpd.\{\alpha,\alpha\}} & {\Grpd.\{\alpha+1,\alpha+1\}}
	\arrow[from=1-1, to=1-2]
	\arrow[from=1-1, to=2-1]
	\arrow[from=1-2, to=2-2]
	\arrow[from=2-1, to=2-2]
    \arrow["\lrcorner"{anchor=center, pos=0.125}, draw=none, from=1-1, to=2-2]
\end{tikzcd}
\]
The functor $\Core : \Cat.\{4,4\} \to \Grpd.\{4,4\}$ is right adjoint to the forgetful functor,
hence it preserves pullbacks,
providing the following pullback square as the image
of the previous
\[
\begin{tikzcd}
	{\Ulow_\alpha} & {\Ulow_{\alpha+1}} \\
	{\U_\alpha} & {\U_{\alpha+1}}
	\arrow[from=1-1, to=1-2]
	\arrow[from=1-1, to=2-1]
    \arrow["\lrcorner"{anchor=center, pos=0.125}, draw=none, from=1-1, to=2-2]
	\arrow[from=1-2, to=2-2]
	\arrow[from=2-1, to=2-2]
\end{tikzcd}
\]
Secondly, we construct a functor $U_\alpha : 1 \to \U_{\alpha+1}$,
which picks out the $(\alpha+1)$-small groupoid $\Core(\Grpd.\{\alpha,\alpha\})$.
It follows that we have the required isomorphism
\[ 1 . U_\alpha \iso \U_\alpha \]
\end{definition}

\begin{definition}[$\Sigma$-types \done{https://github.com/sinhp/HoTTLean/blob/31133dd5b25226ea897f8aa5e2e43b61392459eb/HoTTLean/Groupoids/UHom.lean\#L50}]
For the formation rule of sigma types,
we require for each pair of functors $A : \Gamma \to \U$
and $B : \Gamma . A \to \U$,
a functor $\Sigma_A B : \Gamma \to \U$.
This amounts to constructing for each pair of functors
$A' : \Gamma \to \Grpd$
and $B' : \int A' \to \Grpd$,
a functor $\Sigma_{A'} B' : \Gamma \to \Grpd$.
This can be done ``pointwise'' using Grothendieck constructions.
\[ \textstyle \Sigma_{A'} B' (x) = \int (B |_{A' x} : A' x \to \smallint A' \to \Grpd)\]
One checks that this operation is stable under substitution.
This construction of $\Sigma$-types is automatically universe-polymorphic:
if $A$ is $u$-small and $B$ is $v$-small, then $\Sigma_A B$ is $\max (u,v)$-small.

Rather than constructing unalgebraic operations that interpret the introduction
and elimination rules,
we construct an isomorphism
\[\textstyle\int (B' : \smallint A' \to \Grpd) \iso \int (\Sigma_{A'} B' : \Gamma \to \Grpd)\]
\end{definition}

\begin{definition}[$\Pi$-types \done{https://github.com/sinhp/HoTTLean/blob/31133dd5b25226ea897f8aa5e2e43b61392459eb/HoTTLean/Groupoids/UHom.lean\#L47}]
Our unalgebraic construction of $\Pi$-types
uses the construction of $\Sigma$-types:
the formation of $\Pi$-types is defined by
taking ``pointwise sections'' of the $\Sigma$-types.
As a result of constructing $\Pi$-types this way,
it is easier to not construct $\Pi$-types polymorphically, rather
if $A$ is $u$-small and $B$ is $u$-small, then $\Pi_A B$ is $u$-small.
Then, to obtain universe-polymorphic $\Pi$-types,
we simply combine the above construction with universe lifts.
\end{definition}

\begin{definition}[$\Id$-types \done{https://github.com/sinhp/HoTTLean/blob/31133dd5b25226ea897f8aa5e2e43b61392459eb/HoTTLean/Groupoids/UHom.lean\#L71}]
  \label{def:elementary-id-type}
To construct $\Id$-types in the groupoid model,
it suffices to identify a cylinder structure on the category of groupoids,
and $\Path$-types using the interval (\Cref{thm:id-elim-2}).
The cylinder endofunctor $\cyl : \CC \to \CC$ is constructed
by taking products with the walking isomorphism $I = \{\star \iso \star\}$
\[ \cyl X = X \times I\]
For the formation rule,
let $A : \Gamma \to \Grpd$ be a type and $a, b : \Gamma \to \PGrpd$ be two terms
in $A$.
We construct a functor $\Id_A (a,b) : \Gamma \to \Grpd$ such that
pointwise, it takes the discrete groupoid of morphisms
\[ \Id_A (a,b) (x) = \Hom_{A x} (a x, b x)\]
The $\path$-$\unpath$ bijection follows from
\[ \Hom_{A x} (a x, b x) = \{ f : I \to A x \st f \circ \delta_0 = a x, f \circ \delta_1  = b x\}\]

We construct a normal Hurewicz structure on $\uu : \Ulow \to \U$.
For any groupoid $\Gamma$ and functors
$p : I \times \Gamma \to \U$
and $p_0 : \Gamma \to \Ulow$ such that $\uu \circ p_0 = p \circ (0 \times \Gamma)$,
we construct a lift $l$:
\[ \begin{tikzcd}
	\Gamma & \Ulow \\
	{I \times \Gamma} & \U
	\arrow["{p_0}", from=1-1, to=1-2]
	\arrow["{0 \times \Gamma}"', from=1-1, to=2-1]
	\arrow[from=1-2, to=2-2]
	\arrow["l"{description}, dashed, from=2-1, to=1-2]
	\arrow["p"', from=2-1, to=2-2]
\end{tikzcd}\]
This is provided by the canonical split 
isofibration structure on $\uu : \Ulow \to \U$,
viewing $\Ulow \iso \int i$ as a Grothendieck construction over $\U$.
It follows that these lifts are \emph{normal},
meaning that for any $A : \Gamma \to \U$
and $a : \Gamma \to \Ulow$ such that $\uu \circ a = A$,
the lift provided for the following diagram
\[
\begin{tikzcd}
	\Gamma && \Ulow \\
	{I \times \Gamma} & \Gamma & \U
	\arrow["a", from=1-1, to=1-3]
	\arrow["{0 \times \Gamma}"', from=1-1, to=2-1]
	\arrow[from=1-3, to=2-3]
	\arrow["l"{description}, dashed, from=2-1, to=1-3]
	\arrow["{! \times \Gamma}"', from=2-1, to=2-2]
	\arrow["a"{description}, from=2-2, to=1-3]
	\arrow["A"', from=2-2, to=2-3]
\end{tikzcd}
\]
factors as $l = a \circ (! \times \Gamma)$.
\end{definition}

\section{Algebraic models}\label{sec:StrSem}

Here we provide an alternative algebraic formulation of an MLTT model,
following \cite{awodey2025}, but working in a $\pi$-clan (\Cref{def:pi-preclan}) instead.
For this section, we fix a $\pi$-clan $(\CC, \RR)$.

\begin{definition}[\done{https://github.com/sinhp/HoTTLean/blob/TYPES2026/HoTTLean/Model/Structured/StructuredUniverse.lean\#L20}]
  \label{def:algebraic-universe}
  An \emph{algebraic universe} in $(\CC,\RR)$ is a universe $\uu : \Ulow \to \U$ in
  $\CC$ which is also an $\RR$-map.
  We will just say a universe $\uu$ is $\RR$-algebraic, or algebraic when the $\pi$-clan
  is clear from the context.
\end{definition}

We can replicate the MLTT type formers described in \cite{awodey2025} in this setting,
making use of our general theory of polynomial functors.

\begin{definition}[$\Unit$-types]
  \label{def:alg-unit-type-mltt}
  Suppose $\CC$ has a terminal object $1$.
  We say that an algebraic universe $\uu : \Ulow \to \U$ admits
  \emph{algebraic $\Unit$-types} when we have maps
  $\Unit : 1 \to \Ulow$ and $\unit : 1 \to \U$, such that
  \[\begin{tikzcd}
    1 & \Ulow \\
    1 & \U
    \arrow["\unit", from=1-1, to=1-2]
    \arrow[equals, from=1-1, to=2-1]
    \arrow["\lrcorner"{anchor=center, pos=0.125}, draw=none, from=1-1, to=2-2]
    \arrow["\uu", from=1-2, to=2-2]
    \arrow["\Unit"', from=2-1, to=2-2]
  \end{tikzcd}\]
  is a pullback square.
\end{definition}

% In that case, the relevant pullback square is the following.
% \[
% \begin{tikzcd}
% 	{P_\vv \Ulow} & \Ulow \\
% 	{P_\vv \U} & \U
% 	\arrow["\lam", from=1-1, to=1-2]
% 	\arrow["{P_\vv \vv}"', from=1-1, to=2-1]
% 	\arrow["\lrcorner"{anchor=center, pos=0.125}, draw=none, from=1-1, to=2-2]
% 	\arrow["\vv", from=1-2, to=2-2]
% 	\arrow["\Pi"', from=2-1, to=2-2]
% \end{tikzcd}
% \]
% A more precise correspondence between an elementary $\Pi$-type structure
% and an algebraic one is given in \cite{hua2026}.
Recall that since $\uu$ is an $\RR$-map and $(\CC,\RR)$ is a $\pi$-clan,
we can take $\uu$ to be the signature for a polynomial functor $P_\uu : \CC \to \CC$
(\Cref{def:polynomial-functor-1}).

\begin{definition}[$\Sigma$-types, \done{https://github.com/sinhp/HoTTLean/blob/TYPES2026/HoTTLean/Model/Structured/StructuredUniverse.lean\#L651}]
    \label{def:alg-sig-type-mltt}
    
    Recall the definition of polynomial composition from \Cref{def:pcomp}.
    We consider the polynomial composition of $(\uu : \Ulow \to \U)$ with itself.
    For notational convenience, we write $Q$ for the domain of the composition.
    We say that an
    algebraic universe $\uu:\Ulow\to \U$ admits
    \emph{algebraic $\Sigma$-types} when we have maps
    $\Sigma : P_\uu \U \to \U$ and $\pair: Q \to \Ulow$,
    such that
    \[ \begin{tikzcd}
      Q & \Ulow \\
      {P_\uu\U} & \U
      \arrow["\pair", from=1-1, to=1-2]
      \arrow["{\uu \pcomp \uu}"', from=1-1, to=2-1]
      \arrow["\lrcorner"{anchor=center, pos=0.125}, draw=none, from=1-1, to=2-2]
      \arrow["\uu", from=1-2, to=2-2]
      \arrow["\Sigma"', from=2-1, to=2-2]
    \end{tikzcd} \]
    is a pullback square.
\end{definition}

\begin{definition}[$\Pi$-types, \done{https://github.com/sinhp/HoTTLean/blob/TYPES2026/HoTTLean/Model/Structured/StructuredUniverse.lean\#L359}]
    \label{def:alg-pi-type-mltt}
    We say that an algebraic universe $\uu:\Ulow\to \U$
    admits \emph{algebraic $\Pi$-types} when we have maps
    $\Pi : P_\uu \U \to \U$ and
    $\lam : P_\uu \Ulow\to \Ulow$, such that
    \[ \begin{tikzcd}
      {P_{\uu}{\Ulow}} & \Ulow \\
      {P_{\uu}{\U}} & \U
      \arrow["\lam", from=1-1, to=1-2]
      \arrow["{{P_{\uu}{\uu}}}"', from=1-1, to=2-1]
      \arrow["\lrcorner"{anchor=center, pos=0.125}, draw=none, from=1-1, to=2-2]
      \arrow["\uu", from=1-2, to=2-2]
      \arrow["\Pi"', from=2-1, to=2-2]
    \end{tikzcd} \]
    is a pullback square.
\end{definition}

% Like with elementary $\Pi$-types,
% our formalisation formulates a more general definition
% using three universes $\uu_0 : \Ulow_0 \to \U_0$,
% $\uu_1 : \Ulow_1 \to \U_1$ and $\uu_2 : \Ulow_2 \to \U_2$.
% \[\begin{tikzcd}
% 	{P_{\uu_0}{\Ulow_1}} & {\Ulow_2} \\
% 	{P_{\uu_0}{\U_1}} & {\U_2}
% 	\arrow["\lam", from=1-1, to=1-2]
% 	\arrow["{{P_{\uu_0}{\uu_1}}}"', from=1-1, to=2-1]
% 	\arrow["{\uu_2}", from=1-2, to=2-2]
% 	\arrow["\Pi"', from=2-1, to=2-2]
% \end{tikzcd}\]


% Let us begin by recalling what we mean by a universe having algebraic identity types,
% from \cite{awodey2025}.
% This could be an application of the theory of hom-types from \Cref{sec:cat-as-types},
% where the twisted structure is now degenerate,
% but we repeat the definitions for the sake of clarity.

% \Cref{def:algebraic-id-types} is reformulated in an elementary style in \Cref{def:elementary-id-type}.

\begin{definition}\label{def:algebraic-id-form-intro}
  We will say that an algebraic universe $\uu : \Ulow \to \U$
  admits identity formation and introduction when
  we have a pair of maps $i : \Ulow \to \Ulow$
  and $\Id : \Ulow \times_\U \Ulow \to \Ulow$
  forming a commutative square
  \[ \begin{tikzcd}
    \Ulow & \Ulow \\
    {\Ulow \times_\U \Ulow} & \U
    \arrow["i", from=1-1, to=1-2]
    \arrow["\Delta"', from=1-1, to=2-1]
    \arrow[from=1-2, to=2-2]
    \arrow["\Id"', from=2-1, to=2-2]
  \end{tikzcd} \]
  where the map $\Delta : \Ulow \to {\Ulow \times_\U \Ulow}$ is the diagonal. 
\end{definition}

Assuming that we have structure for formation and introduction of
identity types on a universe $\uu : \Ulow \to \U$,
we can take the pullback of $\Id$ to produce a comparison map
$\rfl : \Ulow \to I$ that factors the diagonal through the pullback.
\[\begin{tikzcd}
  \Ulow && \\
  & I & \Ulow \\
  & {\Ulow \times_\U \Ulow} & \U
  \arrow["\rfl"{description}, dashed, from=1-1, to=2-2]
  \arrow["i", bend left, from=1-1, to=2-3]
  \arrow["\Delta"', bend right, from=1-1, to=3-2]
  \arrow[from=2-2, to=2-3]
  \arrow["d_I"', from=2-2, to=3-2]
  \arrow["\lrcorner"{anchor=center, pos=0.125}, draw=none, from=2-2, to=3-3]
  \arrow[from=2-3, to=3-3]
  \arrow["\Id"', from=3-2, to=3-3]
\end{tikzcd}\]
Let us write $\pp = \uu \circ \pi_0 \circ d_I = \uu \circ \pi_1 \circ d_I : I \to \U$ for the composition 
\[ \begin{tikzcd}
	I & {\Ulow \times_\U \Ulow} & \Ulow & \U
	\arrow["{d_I}", from=1-1, to=1-2]
	\arrow["{\pi_0}", shift left, from=1-2, to=1-3]
	\arrow["{\pi_1}"', shift right, from=1-2, to=1-3]
	\arrow["\uu", from=1-3, to=1-4]
\end{tikzcd}\]
Consider the following commuting triangle.
\[ \begin{tikzcd}
	\Ulow & I \\
	& \U
	\arrow["\rfl", from=1-1, to=1-2]
	\arrow["{\uu}"', from=1-1, to=2-2]
	\arrow["\pp", from=1-2, to=2-2]
\end{tikzcd} \]
Both $(\uu : \Ulow \to \U)$ and $(\pp : I \to \U)$ are
polynomial signatures in $\Poly_{(\CC,\RR)}$,
since $(\CC,\RR)$ is a $\pi$-clan and they are both $\RR$-maps
($\pp$ is a composition of $\RR$-maps).
Using \Cref{def:vert-trans},
this triangle gives rise to a vertical natural transformation
between the polynomial functors.
\[P_{\rfl} : P_{\pp} \to P_{\uu}\]
We consider the naturality square of $P_{\rfl}$
for the morphism $\uu : \Ulow \to \U$.
\[ \begin{tikzcd}
	{ P_{\pp} \Ulow} & {P_{\uu} \Ulow} \\
	{ P_{\pp} \U} & {P_{\uu} \U}
	\arrow["{{P_{\rfl} \Ulow}}", from=1-1, to=1-2]
	\arrow["{{ P_{\pp} \uu}}"', from=1-1, to=2-1]
	\arrow["{{P_{\uu} \uu}}", from=1-2, to=2-2]
	\arrow["{{P_{\rfl} \U}}"', from=2-1, to=2-2]
\end{tikzcd} \]
Since $P_\uu$ preserves $\RR$-maps
(\Cref{prop:polynomial-functor-preclan-morphism}),
the map $P_\uu \uu : P_\uu \Ulow \to P_\uu \U$ is an $\RR$-map,
giving us a pullback $L$:
  \begin{equation}
    \label{eq:id-elimination-2}
    \begin{tikzcd}
    { P_{\pp} \Ulow} && \\
    & L & {P_{\uu} \Ulow} \\
    & { P_{\pp} \U} & {P_{\uu} \U}
    \arrow["p", dashed, from=1-1, to=2-2]
    \arrow["{P_{\rfl} \Ulow}", bend left, from=1-1, to=2-3]
    \arrow["{P_{\pp} \uu}"', bend right, from=1-1, to=3-2]
    \arrow["\phi", from=2-2, to=2-3]
    \arrow["\psi"', from=2-2, to=3-2]
    \arrow["\lrcorner"{anchor=center, pos=0.125}, draw=none, from=2-2, to=3-3]
    \arrow["{P_{\uu} \uu}", from=2-3, to=3-3]
    \arrow["{P_{\rfl} \U}"', from=3-2, to=3-3]
  \end{tikzcd}
  \end{equation}
By the universal property of polynomials
and commutativity of the square,
$\psi$ and $\phi$ are equivalent to the data of three maps
\[\psi_1 = \phi_1 : L \to \U
\quad \psi_2 : \Id_A \to \U
\quad \phi_2 : {A} \to \Ulow\]
where $A := L \opext \psi_1 \to L$ is
given by pullback of $\uu : \Ulow \to \U$ along $\psi_1 : L \to \U$,
and $\Id_A := L \times_U I$ is the pullback in the following diagram
\begin{equation}\label{eq:id-elimination-1}
  \begin{tikzcd}
	\Ulow & {{A}} & \Ulow \\
	\U & {\Id_{A}} & I \\
	& L & \U
	\arrow["\uu"', from=1-1, to=2-1]
	\arrow["{\phi_2}"', from=1-2, to=1-1]
	\arrow[from=1-2, to=1-3]
	\arrow["{\rfl}", from=1-2, to=2-2]
	\arrow["\lrcorner"{anchor=center, pos=0.125}, draw=none, from=1-2, to=2-3]
	\arrow["\rfl", from=1-3, to=2-3]
	\arrow["{\uu}", bend left = 50, from=1-3, to=3-3]
	\arrow["{\psi_2}", from=2-2, to=2-1]
	\arrow[from=2-2, to=2-3]
	\arrow[from=2-2, to=3-2]
	\arrow["\lrcorner"{anchor=center, pos=0.125}, draw=none, from=2-2, to=3-3]
	\arrow["\pp", from=2-3, to=3-3]
	\arrow["{\psi_1}"', from=3-2, to=3-3]
\end{tikzcd}
\end{equation}

\begin{lemma}
  \label{lem:psi-fibration}
  The morphism $\psi_1 : L \to \U$ is an $\RR$-map.
\end{lemma}
\begin{proof}
  The map $\psi_1 : L \to \U$ is a composition of two maps
  $\psi : L \to P_\pp \U$ and
  $\fstProj : P_\pp \U \to \U$,
  the first projection of the polynomial from \Cref{prop:poly-up}.
  The first map $\psi$ is a pullback of $P_\uu \uu$,
  which is an $\RR$-map
  since $P_\uu$ preserves $\RR$-maps
  (\Cref{prop:polynomial-functor-preclan-morphism})
  As for $\fstProj$,
  recall that evaluating a polynomial at the terminal object
  recovers the base of the signature
  $P_\pp 1 \iso \U$.
  The map $\fstProj : P_\pp \U \to \U$ is therefore isomorphic to
  $P_\pp ! : P_\pp \U \to P_\pp 1$,
  where $! : \U \to 1$.
  Finally, $P_\pp !$ is an $\RR$-map because $P_\pp$ preserves $\RR$-maps
  and $!$ is an $\RR$-map.
\end{proof}

\begin{proposition}[Identity elimination]\label{prop:id-elim-equiv}
  Suppose $\uu : \Ulow \to \U$ is an algebraic universe that admits
  identity formation and introduction.
  The following are equivalent (structures)
  % for a universe $\uu : \Ulow \to \U$ that admits
  % identity formation and introduction.
  \begin{itemize}
    \item A section $j : L \to P_{\pp} \Ulow$
    of the comparison map $p : P_{\pp} \Ulow \to L$
    from \cref{eq:id-elimination-2}.
    \[ p \circ j = \id\]
    \item A diagonal lift $j : \Id_{A} \to \Ulow$
    as indicated in the following diagram,
    extracted from \cref{eq:id-elimination-1}.
    \[\begin{tikzcd}
        {{A}} & \Ulow \\
        {\Id_A} & \U
        \arrow["{{\phi_2}}", from=1-1, to=1-2]
        \arrow["{{\rfl}}"', from=1-1, to=2-1]
        \arrow["\uu", from=1-2, to=2-2]
        \arrow["j"{description}, dashed, from=2-1, to=1-2]
        \arrow["{{\psi_2}}"', from=2-1, to=2-2]
      \end{tikzcd}\]
    \item A diagonal lift $j' : \Id_{A} \to C$
    as indicated in the following diagram.
    \[\begin{tikzcd}
        {{A}} & C & \Ulow \\
        {\Id_{A}} & {\Id_{A}} & \U
        \arrow["r", dashed, from=1-1, to=1-2]
        \arrow["{{{\phi_2}}}", bend left, from=1-1, to=1-3]
        \arrow["{{{\rfl}}}"', from=1-1, to=2-1]
        \arrow[from=1-2, to=1-3]
        \arrow[from=1-2, to=2-2]
        \arrow["\lrcorner"{anchor=center, pos=0.125}, draw=none, from=1-2, to=2-3]
        \arrow["\uu", from=1-3, to=2-3]
        \arrow["{j'}"{description}, dashed, from=2-1, to=1-2]
        \arrow[equals, from=2-1, to=2-2]
        \arrow["{{{\psi_2}}}"', from=2-2, to=2-3]
      \end{tikzcd}\]
  \end{itemize}
\end{proposition}

\begin{definition}[$\Id$-types]\label{def:algebraic-id-types}
  If an algebraic universe admits identity formation and introduction and
  any one of the equivalent conditions of \Cref{prop:id-elim-equiv} is met (elimination),
  we will say that the universe has $\Id$-type structure.
\end{definition}

\begin{proposition}
  If $\uu : \Ulow \to \U$ is an $\RR$-algebraic universe,
  and $\uu$ admits algebraic $\Unit$-types
  (respectively $\Sigma$-types, $\Pi$-types, $\Id$-types),
  it also admits unalgebraic $\Unit$-types
  (respectively $\Sigma$-types, $\Pi$-types, $\Id$-types).
\end{proposition}
\begin{proof}
  This is routine (cf. \cite{awodey2025, hua2026}).
\end{proof}
As a result, when the universe is algebraic,
we can drop the adjective ``algebraic'' and ``unalgebraic''
for type formers.

\subsection*{Constructing algebraic models}

\begin{definition}\label{def:alg-universe-classifying-map}
  Suppose $\uu : \Ulow \to \U$ is a universe.
  We denote the set of pullbacks of $\uu$ by $\RR_\U$
  and call $\RR_\U$-maps \emph{small fibrations}.
  (If $\uu$ is $\RR$-algebraic then $\RR_\U \subseteq \RR$.)
  \[
    \RR_\U := \{ A : X \to \Gamma \st \begin{tikzcd}
    X & \Ulow \\
    \Gamma & \U
    \arrow["\exists", dashed, from=1-1, to=1-2]
    \arrow["A"', from=1-1, to=2-1]
    \arrow["\lrcorner"{anchor=center, pos=0.125}, draw=none, from=1-1, to=2-2]
    \arrow["\uu", from=1-2, to=2-2]
    \arrow["\exists"', dashed, from=2-1, to=2-2]
  \end{tikzcd}\}
  \]
  Then we say that $\uu$ \emph{has chosen classifying maps}
  when every small fibration $A : X \to \Gamma$ in $\RR_\U$
  is equipped with a {chosen classifying map}
  $\ulcorner A \urcorner : \Gamma \to \U$,
  so that
  \[ \begin{tikzcd}
    X & \Ulow \\
    \Gamma & \U
    \arrow[from=1-1, to=1-2]
    \arrow["A"', from=1-1, to=2-1]
    \arrow["\lrcorner"{anchor=center, pos=0.125}, draw=none, from=1-1, to=2-2]
    \arrow["\uu", from=1-2, to=2-2]
    \arrow["{{\ulcorner A \urcorner}}"', from=2-1, to=2-2]
  \end{tikzcd} \]
  is a pullback square.
\end{definition}
When the axiom of choice is available,
this condition is trivial.
In many models, such as the groupoid model with $\uu : \Ulow \to \U$
as the universe of small isofibrations,
it is not constructively true that the universe has chosen classifying maps.
In this case, we should really work in a setting where $\RR_\U$
is not just a class of maps satisfying a property,
but a class of maps with structure.
% Corresponding to this,
% the structure of a small split isofibration is preserved under 
% pullback, composition, isomorphism, and pushforward.
For now, we will make it clear where chosen classifying maps are being used,
and have the groupoid model as a non-constructive example.
% For the universe of small split isofibrations in the category of groupoids
% (\Cref{sec:universe-in-groupoid-model}),
% this condition is constructively satisfied by taking the fibres functor
% of a split isofibration (cf. \Cref{prop:universal-opfibration-classifies}).

\begin{proposition}\label{prop:constructing-alg-models}
  Suppose $(\CC,\RR)$ is a $\pi$-clan and $\uu : \Ulow \to \U$ is an $\RR$-algebraic universe
  that has chosen classifying maps.
  If $\RR_\U$ is a $\pi$-preclan then
  $\uu$ admits $\Unit$-types, $\Sigma$-types and $\Pi$-types.
\end{proposition}
\begin{proof}
  Since $\uu$ has chosen classifying maps,
  to show that $\uu$ admits $\Unit$-types,
  we only need to note that the identity $1 \to 1$
  is a small fibration.

  To show that $\uu$ admits $\Sigma$-types,
  we can simply recall that $\uu \pcomp \uu$
  is defined as a composition of pullbacks of small fibrations,
  and is therefore also a small fibration.

  To show that $\uu$ admits $\Pi$-types,
  it suffices to show that $P_\uu \uu$ is a small fibration.
  This follows from \Cref{lem:polynomial-functor-preserves-maps} with $\SS = \RR_\U$,
  which says that $P_\uu : (\RR(1),\RR_\U) \to (\RR(1),\RR_\U)$ is a preclan morphism.
\end{proof}

Modulo concerns about constructivism,
\Cref{prop:constructing-alg-models} makes models of MLTT much easier to set up.
For example,
when working with the groupoid model,
one only needs to check that the class of small split isofibrations
(meaning the maps $F : \CC \to \DD$ on which there merely exists the
structure of a small split isofibration $(F,l)$)
are closed under isomorphisms, compositions, and pushforwards.
(A stronger claim is also true: the structure of small split isofibrations
are preserved under these operations.)
The stability of type formers under substitution automatically
follows from the general theory of algebraic models.
% To construct algebraic $\Id$-types,
% we will use algebraic $\Path$-types,
% which we present in \Cref{sec:path-types}.
For constructing algebraic $\Id$-types,
one could either consider the usual factorisation conditions,
which we do below,
or algebraic $\Path$-types, as we present in \Cref{sec:path-types}.

\begin{definition}[Anodyne maps]
  Let $\RR$ be a class of maps in a category $\CC$.
  A morphism $f : X \to Y$ is $\RR$-anodyne
  when $f$ has the left lifting property against all $\RR$-maps,
  meaning that for any $\RR$-map $d : B \to A$, and morphisms
  $b : X \to B$ and $a : Y \to A$ such that $d \circ b = a \circ f$,
  there is a diagonal lift $t : Y \to B$.
  We draw $\RR$-maps with a double head $\twoheadrightarrow$
  and anodyne maps with a tail $\rightarrowtail$. 
  \[
    \begin{tikzcd}
    X & B \\
    Y & A
    \arrow[from=1-1, to=1-2]
    \arrow[tail, from=1-1, to=2-1]
    \arrow[two heads, from=1-2, to=2-2]
    \arrow[dashed, from=2-1, to=1-2]
    \arrow[from=2-1, to=2-2]
  \end{tikzcd}
  \]
  % A pretribe $(\CC,\RR)$ is a preclan such that 
  % \begin{itemize}
  %   \item every map $f : X \to Y$ between $\RR$-objects factors as an
  %   anodyne map followed by an $\RR$-map and
  %   \item anodyne maps are stable under pullback along $\RR$-maps.
  % \end{itemize}
\end{definition}

% \begin{definition}
%   A subtribe $(\CC,\RR,\RR_\U)$ consists of 
%   \begin{itemize}
%     \item a clan $(\CC,\RR)$ such that $\RR$-anodyne maps are stable under pullback
%     \item a preclan $\RR_\U \subseteq \RR$,
%     \item 
%   \end{itemize}
%   Let us say that $\RR_U \subseteq \RR$ is a pretribe in $\RR$
%   when any 
% \end{definition}

\begin{lemma}\label{lem:anodyne-pullback}
  If $(\CC,\RR)$ is a $\pi$-clan,
  then anodyne maps are stable under pullback along $\RR$-maps.
\end{lemma}
\begin{proof}
  See \cite[Lemma 29]{awodey2018}.
\end{proof}

% \begin{proposition}\label{prop:pi-pretribe-equiv}
%   Suppose $(\CC,\RR)$ is a $\pi$-clan and $\uu : \Ulow \to \U$
%   is an $\RR$-algebraic universe.
%   Suppose that $\RR_\U$, the class of pullbacks of $\uu$, forms a preclan.
%   Then the following conditions are equivalent
%   \begin{enumerate}
%     % \item The diagonal $\Delta : \Ulow \to \Ulow \times_\U \Ulow$
%     %   factors as an $\RR$-anodyne map
%     %   followed by an $\RR_\U$-map.
%     \item For any $\RR_\U$-map $A \to X$,
%       the diagonal $\Delta : A \to A \times_X A$ factors as an $\RR$-anodyne map
%       followed by an $\RR_\U$-map.
%     % \item Given objects $X$ in $\CC$,
%     %   $A$ in $\RR_\U (X)$, $B$ in $\CC / X$ and a map $f : B \to A$ in $\CC / X$,
%     \item Given an object $X$ in $\CC$
%       and a morphism $f : B \to A$ in $\RR_\U (X)$,
%       $f$ factors as an $\RR$-anodyne map followed by an $\RR_\U$-map
%       \[
%         \begin{tikzcd}
%           B && A \\
%           & {B'}
%           \arrow["f", from=1-1, to=1-3]
%           \arrow[tail, from=1-1, to=2-2]
%           \arrow["{\in \RR_\U}"', two heads, from=2-2, to=1-3]
%         \end{tikzcd}
%       \]
%   \end{enumerate}
% \end{proposition}
% \begin{proof}
%   (2) $\Rightarrow$ (1) is straightforward,
%   by viewing $\Delta : A \to A \times_X A$ as a morphism in $\RR_\U (X)$.
  
%   % by viewing $\Delta : \Ulow \to \Ulow \times_\U \Ulow$ as a morphism in $\RR_\U (\U)$.
  
%   % (1) $\Rightarrow$ (2).
%   % We use a similar construction to that of \cite[Lemma 11]{gambino2008}.
%   % Let us denote the given factorisation of $\Delta$ as follows.
%   % \[ \begin{tikzcd}
%   %     \Ulow && {\Ulow \times_\U \Ulow} \\
%   %     & I
%   %     \arrow["\Delta", from=1-1, to=1-3]
%   %     \arrow["\rfl"', tail, from=1-1, to=2-2]
%   %     \arrow["{(\src,\trg)}"', two heads, from=2-2, to=1-3]
%   %   \end{tikzcd} \]
%   % Let $A$ be an object in $\RR_\U(X)$.
%   % Then there is a classifying map $X \to \U$ for $A$.
%   % \[\begin{tikzcd}
%   %   A & \Ulow \\
%   %   X & \U
%   %   \arrow["t", dashed, from=1-1, to=1-2]
%   %   \arrow[from=1-1, to=2-1]
%   %   \arrow["\lrcorner"{anchor=center, pos=0.125}, draw=none, from=1-1, to=2-2]
%   %   \arrow["\uu", from=1-2, to=2-2]
%   %   \arrow["{\ulcorner A \urcorner}"', dashed, from=2-1, to=2-2]
%   % \end{tikzcd}\]
%   % By \Cref{lem:anodyne-pullback},
%   % the factorisation on $\Delta$ pulls back to one on the diagonal of $A$ (over $X$).
%   % \[ \begin{tikzcd}
%   %   A & \Ulow \\
%   %   {\Id_A} & I \\
%   %   {A \times_X A} & {\Ulow \times_\U \Ulow}
%   %   \arrow["t", from=1-1, to=1-2]
%   %   \arrow["\rfl"', tail, from=1-1, to=2-1]
%   %   \arrow["\lrcorner"{anchor=center, pos=0.125}, draw=none, from=1-1, to=2-2]
%   %   \arrow["\rfl", tail, from=1-2, to=2-2]
%   %   \arrow[from=2-1, to=2-2]
%   %   \arrow["{{(\src,\trg)}}"', two heads, from=2-1, to=3-1]
%   %   \arrow["\lrcorner"{anchor=center, pos=0.125}, draw=none, from=2-1, to=3-2]
%   %   \arrow["{{(\src,\trg)}}", two heads, from=2-2, to=3-2]
%   %   \arrow["{{(t,t)}}"', from=3-1, to=3-2]
%   % \end{tikzcd}\]
%   % Let $f : B \to A$ be a map in $\RR_\U(X)$.
%   % Consider the following two pullbacks.
%   % \[\begin{tikzcd}
%   %     B && A \\
%   %     {B \times_A \Id_A} && {\Id_A} \\
%   %     B && A
%   %     \arrow["f", from=1-1, to=1-3]
%   %     \arrow["r"', tail, from=1-1, to=2-1]
%   %     \arrow["\lrcorner"{anchor=center, pos=0.125}, draw=none, from=1-1, to=2-3]
%   %     \arrow["\rfl", from=1-3, to=2-3]
%   %     \arrow["{f'}"', from=2-1, to=2-3]
%   %     \arrow[from=2-1, to=3-1]
%   %     \arrow["\lrcorner"{anchor=center, pos=0.125}, draw=none, from=2-1, to=3-3]
%   %     \arrow["\src", from=2-3, to=3-3]
%   %     \arrow["f"', from=3-1, to=3-3]
%   %   \end{tikzcd}\]  
%   % We factorise $f$ as
%   % \[
%   % \begin{tikzcd}
%   %   B && A \\
%   %   & {B \times_A \Id_A}
%   %   \arrow["f", from=1-1, to=1-3]
%   %   \arrow["r"', tail, from=1-1, to=2-2]
%   %   \arrow["{\trg \circ f'}"', two heads, from=2-2, to=1-3]
%   % \end{tikzcd}
%   % \]
  
%   (1) $\Rightarrow$ (2).
%   We use a similar construction to that of \cite[Lemma 11]{gambino2008}.
%   Let $f : B \to A$ be a map in $\RR_\U(X)$.
%   Let us denote the given factorisation of the diagonal of $A$ as follows.
%   \[ \begin{tikzcd}
%       A && {A \times_X A} \\
%       & \Id_A
%       \arrow["\Delta", from=1-1, to=1-3]
%       \arrow["\rfl"', tail, from=1-1, to=2-2]
%       \arrow["{(\src,\trg)}"', two heads, from=2-2, to=1-3]
%     \end{tikzcd} \]
%   Consider the following two pullbacks.
%   \[\begin{tikzcd}
%       B && A \\
%       {B \times_A \Id_A} && {\Id_A} \\
%       B && A
%       \arrow["f", from=1-1, to=1-3]
%       \arrow["r"', tail, from=1-1, to=2-1]
%       \arrow["\lrcorner"{anchor=center, pos=0.125}, draw=none, from=1-1, to=2-3]
%       \arrow["\rfl", from=1-3, to=2-3]
%       \arrow["{f'}"', from=2-1, to=2-3]
%       \arrow[from=2-1, to=3-1]
%       \arrow["\lrcorner"{anchor=center, pos=0.125}, draw=none, from=2-1, to=3-3]
%       \arrow["\src", from=2-3, to=3-3]
%       \arrow["f"', from=3-1, to=3-3]
%     \end{tikzcd}\]  
%   We factorise $f$ as
%   \[
%   \begin{tikzcd}
%     B && A \\
%     & {B \times_A \Id_A}
%     \arrow["f", from=1-1, to=1-3]
%     \arrow["r"', tail, from=1-1, to=2-2]
%     \arrow["{\trg \circ f'}"', two heads, from=2-2, to=1-3]
%   \end{tikzcd}
%   \]
% \end{proof}

\begin{proposition}\label{prop:constructing-alg-id-types}
  Suppose $(\CC,\RR)$ is a $\pi$-clan and
  $\uu : \Ulow \to \U$ is an $\RR$-algebraic universe that has chosen classifying maps.
  If the diagonal on the universe
  $\Delta : \Ulow \to \Ulow \times_\U \Ulow$
  factors as an $\RR$-anodyne\footnote{
    For the sake of avoiding the axiom of choice,
    we also need to assume that anodyne maps
    come equipped with certain chosen lifts.
    Let us do so implicitly.} %TODO
  map followed by an $\RR_\U$-map
  then $\uu$ admits $\Id$-types.
\end{proposition}
\begin{proof}
  Let us denote the given factorisation of $\Delta$ as follows.
  \[ \begin{tikzcd}
      \Ulow && {\Ulow \times_\U \Ulow} \\
      & I
      \arrow["\Delta", from=1-1, to=1-3]
      \arrow["\rfl"', tail, from=1-1, to=2-2]
      \arrow["{(\src,\trg)}"', two heads, from=2-2, to=1-3]
    \end{tikzcd} \]
  Identity formation and introduction follows from
  $(\src,\trg) : I \to {\Ulow \times_\U \Ulow}$
  being in $\RR_\U$.
  \[ \begin{tikzcd}
    \Ulow & I & \Ulow \\
    & {{\Ulow \times_\U \Ulow}} & \U
    \arrow["\rfl", from=1-1, to=1-2]
    \arrow["\Delta"', from=1-1, to=2-2]
    \arrow[dashed, from=1-2, to=1-3]
    \arrow[from=1-2, to=2-2]
    \arrow[from=1-3, to=2-3]
    \arrow["{\ulcorner I \urcorner}"', dashed, from=2-2, to=2-3]
    \arrow["\lrcorner"{anchor=center, pos=0.125}, draw=none, from=1-2, to=2-3]
  \end{tikzcd}\]
  By the second (or third) condition in \Cref{prop:id-elim-equiv},
  to provide identity elimination,
  it suffices to show that the map $A \to \Id_A$
  from the ``universal elimination'' diagram
  \labelcref{eq:id-elimination-1}, reproduced below, is anodyne.
  \[
    \begin{tikzcd}
    \Ulow & {{A}} & \Ulow \\
    \U & {\Id_{A}} & I \\
    & L & \U
    \arrow["\uu"', from=1-1, to=2-1]
    \arrow["{\phi_2}"', from=1-2, to=1-1]
    \arrow[from=1-2, to=1-3]
    \arrow["{\rfl}", from=1-2, to=2-2]
    \arrow["\lrcorner"{anchor=center, pos=0.125}, draw=none, from=1-2, to=2-3]
    \arrow["\rfl", from=1-3, to=2-3]
    \arrow["{\uu}", bend left = 50, from=1-3, to=3-3]
    \arrow["{\psi_2}", from=2-2, to=2-1]
    \arrow[from=2-2, to=2-3]
    \arrow[from=2-2, to=3-2]
    \arrow["\lrcorner"{anchor=center, pos=0.125}, draw=none, from=2-2, to=3-3]
    \arrow["\pp", from=2-3, to=3-3]
    \arrow["{\psi_1}"', from=3-2, to=3-3]
  \end{tikzcd}\]
  The map $A \to \Id_A$ is (by definition) a pullback of the map
  $\rfl : \Ulow \to I$ along $\psi_1 : L \to \U$.
  By \Cref{lem:anodyne-pullback},
  any pullback of an anodyne map along an $\RR$-map is anodyne.
  As required,
  $\rfl : \Ulow \to I$ is anodyne
  and $\psi_1 : L \to \U$ is an $\RR$-map \Cref{lem:psi-fibration}.
\end{proof}

% We stated the result of \Cref{prop:constructing-alg-id-types}
% propositionally,
% but we could also make the same constructive construction:
% If additionally, the universe has chosen classifying maps,
% lifts for anodyne maps
% and the anodyne-$\RR$-map factorisations are given as structure
% (rather than mere existence),
% then the universe admits $\Id$-types.

% The equivalent condition of \Cref{prop:pi-pretribe-equiv}
% are similar to the condition of being a sub-tribe in \cite{joyal2017},
% except, using the terminology of \cite{joyal2017},
% we require that $(\CC,\RR)$ is a $\pi$-clan instead of a tribe.

\section{Future work}
\label{sec:algebraic-vs-elementary}

\subsection*{HoTT0}
We have not shown that the groupoid model of MLTT is a model of HoTT0,
a fragment of HoTT in which univalence holds only on the set-truncated types \cite{hua2025}. 
The two axioms of HoTT0 are set-truncated univalence and function extensionality.
Though a brute-force proof of the axioms should be fairly straightforward,
we would prefer to first create a general user interface
for this kind of syntax-semantics reasoning,
using the tools developed in \cite{nawrocki2026}.

\subsection*{Algebraic models}
Constructing the groupoid model as an unalgebraic model comes with several problems. 
The main problems are those of type equalities and nested dependent equalities,
which ultimately add several minutes of compile time to the files defining type formers in the model.
Another issue is the code's reusability.
The unalgebraic constructions result in long files of code that pertain only to groupoids.
The unalgebraic approach also requires explicit proofs of
stability of all operations under substitution,
though these equations are usually very easy to prove.

Nonetheless, as a target for writing an interpretation function from syntax to semantics,
the unalgebraic model is ideal.
This is because each judgment in the syntax corresponds closely to a rule
in the unalgebraic semantics.
The HoTTLean library would benefit from having
a translation between the two descriptions of models,
taking advantage of both formulations.

To construct an algebraic groupoid model instead,
we would need to formalise the following.
\begin{enumerate}
  \item The theory of algebraic models, including the theory of polynomial functors in preclans.
  \item The $\pi$-clan of isofibrations (or split isofibrations) in groupoids.
  \item The universe of split small isofibrations.
  \item Small split isofibrations form a $\pi$-preclan.
  \item The pathobject factorisation of the diagonal of an isofibration
    (analogous to \Cref{prop:src-trg-discrete-opfibration}).
  \item Either showing that isofibrations are Hurewicz
    or that the pathobject factorisation is part of a weak factorisation system
    (analogous to \Cref{thm:lari-opfib-awfs}).
  \item An algebraic-to-unalgebraic model translation.
\end{enumerate}

We have strong reasons to believe that all of the items
in the list admit quite straightforward and ``nice'' formalisations
(some of the items have already been formalised),
possibly with the exception of showing that
split isofibrations are closed under pushforwards (as part of (4)).
Since the closure conditions of (4) are only required up to isomorphism
and not equality, we expect that there will be no equalities of types 
when working with groupoids.
Instead, those dependent equalities that are inherent to a strict model of
type theory are automatically and uniformly generated in (7),
resulting in a much faster compile time.
Moreover, every item in this list would make for a good contribution to Mathlib,
and are of general interest to the formalisation community.

Finally, there is the issue of using the axiom of choice when applying
\Cref{prop:constructing-alg-id-types,prop:constructing-alg-models}
to construct the groupoid model.
The propositions are currently stated using the class of maps $\RR_\U$
that are (propositionally) pullbacks of the universe,
meaning that the axiom of choice is required to extract a classifying map.
In Mathlib, the convention is usually to apply choice whenever needed,
but a constructive alternative formulation of the result should also be possible.
We leave this for future work.

\subsection*{Algebraic categorical models}

In \Cref{sec:cat-as-types},
we will develop a blueprint
for extending the HoTTLean formalisation of the groupoid model
to the category of categories $\Cat$.
% The development of \Cref{sec:cat-as-types} is designed in part as a blueprint
% for extending the HoTTLean formalisation of the groupoid model
% to the category of categories $\Cat$.
Indeed, many parts of the groupoid model deserve to be factored through
more general theorems in $\Cat$.
Such results would include the following:
\begin{itemize}
  \item The clans of (groupoidal or discrete) (split) (op)fibrations (\Cref{prop:preclans-in-cat})
    and the clan of isofibrations in $\Cat$.
  \item Straightening and unstraightening of split (op)fibrations (\Cref{thm:opfibration-main}).
  \item The universes of split (op)fibrations (\Cref{def:universal-small-split-opfibration}).
  \item The $\tau$-clan of opfibrations and fibrations in $\Cat$ (\Cref{prop:pushforward-in-cat}).
  \item The (lari, split opfibration) algebraic weak factorisation system (\Cref{thm:lari-opfib-awfs}).
\end{itemize}
Again, these results would make for a more modular library,
and could ultimately be part of a library for automated internal reasoning in $\Cat$.